\documentclass{amsart}
% For dealing with references we use the comment environment
\usepackage{verbatim}
\newenvironment{reference}{\comment}{\endcomment}
%\newenvironment{reference}{}{}
\newenvironment{slogan}{\comment}{\endcomment}
\newenvironment{history}{\comment}{\endcomment}

% For commutative diagrams we use Xy-pic
\usepackage[all]{xy}

% We use 2cell for 2-commutative diagrams.
\xyoption{2cell}
\UseAllTwocells

% We use multicol for the list of chapters between chapters
\usepackage{multicol}

% This is generally recommended for better output
\usepackage{lmodern}
\usepackage[T1]{fontenc}

% For cross-file-references

% Package for hypertext links:
\usepackage{hyperref}

% For any local file, say "hello.tex" you want to link to please

% Theorem environments.
%
\theoremstyle{plain}
\newtheorem{theorem}[subsection]{Theorem}
\newtheorem{proposition}[subsection]{Proposition}
\newtheorem{lemma}[subsection]{Lemma}

\theoremstyle{definition}
\newtheorem{definition}[subsection]{Definition}
\newtheorem{example}[subsection]{Example}
\newtheorem{exercise}[subsection]{Exercise}
\newtheorem{situation}[subsection]{Situation}

\theoremstyle{remark}
\newtheorem{remark}[subsection]{Remark}
\newtheorem{remarks}[subsection]{Remarks}

\numberwithin{equation}{subsection}

% Macros
%
\def\lim{\mathop{\mathrm{lim}}\nolimits}
\def\colim{\mathop{\mathrm{colim}}\nolimits}
\def\Spec{\mathop{\mathrm{Spec}}}
\def\Hom{\mathop{\mathrm{Hom}}\nolimits}
\def\Ext{\mathop{\mathrm{Ext}}\nolimits}
\def\SheafHom{\mathop{\mathcal{H}\!\mathit{om}}\nolimits}
\def\SheafExt{\mathop{\mathcal{E}\!\mathit{xt}}\nolimits}
\def\Sch{\mathit{Sch}}
\def\Mor{\mathop{\mathrm{Mor}}\nolimits}
\def\Ob{\mathop{\mathrm{Ob}}\nolimits}
\def\Sh{\mathop{\mathit{Sh}}\nolimits}
\def\NL{\mathop{N\!L}\nolimits}
\def\CH{\mathop{\mathrm{CH}}\nolimits}
\def\proetale{{pro\text{-}\acute{e}tale}}
\def\etale{{\acute{e}tale}}
\def\QCoh{\mathit{QCoh}}
\def\Ker{\mathop{\mathrm{Ker}}}
\def\Im{\mathop{\mathrm{Im}}}
\def\Coker{\mathop{\mathrm{Coker}}}
\def\Coim{\mathop{\mathrm{Coim}}}

% Boxtimes
%
\DeclareMathSymbol{\boxtimes}{\mathbin}{AMSa}{"02}

%
% Macros for moduli stacks/spaces
%
\def\QCohstack{\mathcal{QC}\!\mathit{oh}}
\def\Cohstack{\mathcal{C}\!\mathit{oh}}
\def\Spacesstack{\mathcal{S}\!\mathit{paces}}
\def\Quotfunctor{\mathrm{Quot}}
\def\Hilbfunctor{\mathrm{Hilb}}
\def\Curvesstack{\mathcal{C}\!\mathit{urves}}
\def\Polarizedstack{\mathcal{P}\!\mathit{olarized}}
\def\Complexesstack{\mathcal{C}\!\mathit{omplexes}}
% \Pic is the operator that assigns to X its picard group, usage \Pic(X)
% \Picardstack_{X/B} denotes the Picard stack of X over B
% \Picardfunctor_{X/B} denotes the Picard functor of X over B
\def\Pic{\mathop{\mathrm{Pic}}\nolimits}
\def\Picardstack{\mathcal{P}\!\mathit{ic}}
\def\Picardfunctor{\mathrm{Pic}}
\def\Deformationcategory{\mathcal{D}\!\mathit{ef}}

\setlength{\emergencystretch}{3em}
\begin{document}
\title[Stacks Project, Tag 07H8]{329 --- Existence of relative divided power envelopes for arbitrary ring-ideal maps}
\maketitle
\noindent Copyright (C) 2005--2025 Johan de Jong. Stacks Project authors. Source: \href{https://stacks.math.columbia.edu/tag/07H8}{Tag 07H8}.

\noindent Editorial difficulty note (not author text). Difficulty H1: Representability requires a solution-set construction, not simply a formal appeal to a free polynomial algebra. The proof constructs a bounded divided-power subalgebra containing all specified images and operations, then verifies compatibility with arbitrary limits to produce the universal envelope..

\noindent Editorial provenance note (not author text). History: excerpt from revision \texttt{a0444\allowbreak 6e57e\allowbreak c1fbc\allowbreak 252a8\allowbreak 71afc\allowbreak ec775\allowbreak 2fb28\allowbreak 07b14}, crystalline.tex, lines 53--108. Reference presentation was adapted; original-author context and core support blocks were retained or added. The mathematical wording is unchanged. Licensed under GNU FDL 1.2 or later; no invariant sections or cover texts. See ../licenses/COPYING. Context blocks reproduce author definitions and supporting results; they are not additional counted problems.

\begin{lemma}
\label{lemma-divided-power-envelope}
Let $(A, I, \gamma)$ be a divided power ring.
Let $A \to B$ be a ring map. Let $J \subset B$ be an ideal
with $IB \subset J$. There exists a homomorphism of
divided power rings
$$
(A, I, \gamma) \longrightarrow (D, \bar J, \bar \gamma)
$$
such that
$$
\Hom_{(A, I, \gamma)}((D, \bar J, \bar \gamma), (C, K, \delta)) =
\Hom_{(A, I)}((B, J), (C, K))
$$
functorially in the divided power algebra $(C, K, \delta)$ over
$(A, I, \gamma)$. Here the LHS is morphisms of divided
power rings over $(A, I, \gamma)$ and the RHS is morphisms of
(ring, ideal) pairs over $(A, I)$.
\end{lemma}

\begin{proof}
Denote $\mathcal{C}$ the category of divided power rings
$(C, K, \delta)$. Consider the functor
$F : \mathcal{C} \longrightarrow \textit{Sets}$ defined by
$$
F(C, K, \delta) =
\left\{
(\varphi, \psi)
\middle|
\begin{matrix}
\varphi : (A, I, \gamma) \to (C, K, \delta)
\text{ homomorphism of divided power rings} \\
\psi : (B, J) \to (C, K)\text{ an }
A\text{-algebra homomorphism with }\psi(J) \subset K
\end{matrix}
\right\}
$$
We will show that
Divided Power Algebra, Lemma \href{https://stacks.math.columbia.edu/tag/07GW}{lemma a version of brown (Tag 07GW)}
applies to this functor which will
prove the lemma. Suppose that $(\varphi, \psi) \in F(C, K, \delta)$.
Let $C' \subset C$ be the subring generated by $\varphi(A)$,
$\psi(B)$, and $\delta_n(\psi(f))$ for all $f \in J$.
Let $K' \subset K \cap C'$ be the ideal of $C'$ generated by
$\varphi(I)$ and $\delta_n(\psi(f))$ for $f \in J$.
Then $(C', K', \delta|_{K'})$ is a divided power ring and
$C'$ has cardinality bounded by the cardinal
$\kappa = |A| \otimes |B|^{\aleph_0}$.
Moreover, $\varphi$ factors as $A \to C' \to C$ and $\psi$ factors
as $B \to C' \to C$. This proves assumption (1) of
Divided Power Algebra, Lemma \href{https://stacks.math.columbia.edu/tag/07GW}{lemma a version of brown (Tag 07GW)}
holds. Assumption (2) is clear
as limits in the category of divided power rings commute with the
forgetful functor $(C, K, \delta) \mapsto (C, K)$, see
Divided Power Algebra, Lemma \href{https://stacks.math.columbia.edu/tag/07GV}{lemma limits (Tag 07GV)} and its proof.
\end{proof}
\end{document}
